\section{Conjuntos e índices}

\begin{longtable}{p{0.20\linewidth}p{0.68\linewidth}}
\toprule
Símbolo & Definición \\
\midrule
\endhead
$\Train$ & conjunto de 138 parcelas con rendimiento observado. \\
$\Target$ & conjunto de 59 parcelas objetivo con rendimiento oculto. \\
$\Target^\star$ & conjunto pseudo-target ocultado dentro de las 138 etiquetas. \\
$\Train^\star$ & etiquetas visibles de una pseudo-competición,
$\Train\setminus\Target^\star$. \\
$x_i$ & vector de covariables de la parcela $i$. \\
$y_i$ & rendimiento observado de parcela $i$, en toneladas por hectárea. \\
$X_{\Train}$ & matriz de covariables de parcelas etiquetadas. \\
$X_{\Target}$ & matriz observable de covariables de las 59 parcelas objetivo. \\
$W$ o $A$ & matriz de pesos/adyacencia de un grafo. \\
$D$ & matriz diagonal de grados del grafo. \\
$L=D-A$ & Laplaciano no normalizado. \\
$d_{ij}$ & distancia entre parcelas $i,j$ en una topología especificada. \\
$\mathcal N_k(q)$ & $k$ vecinos etiquetados seleccionados para query $q$. \\
\bottomrule
\end{longtable}

\section{Métricas de regresión}

Para observaciones $y_i$ y predicciones $\hat y_i$:
\begin{align}
\MSE
&=
\frac{1}{n}\sum_{i=1}^n(y_i-\hat y_i)^2,\\
\RMSE
&=
\sqrt{\MSE},\\
\MAE
&=
\frac{1}{n}\sum_{i=1}^n|y_i-\hat y_i|,\\
R^2
&=
1-
\frac{\sum_i(y_i-\hat y_i)^2}
{\sum_i(y_i-\bar y)^2}.
\end{align}

En reportes con múltiples splits se distinguen:
\begin{align}
\overline{\RMSE}_{split}
&=
\frac{1}{S}\sum_{s=1}^S\RMSE_s,\\
\RMSE_{\mathrm{pooled}}
&=
\sqrt{
\frac{\sum_s\sum_{i\in s}(y_i-\hat y_i)^2}
{\sum_s n_s}
}.
\end{align}
El primer término da igual peso a cada split; el segundo da igual peso a cada fila
pseudo-predicha. Ambos se conservan porque los tamaños de folds grouped pueden ser
desiguales.

\section{Spearman}

Para pares de variables $u,v$, Spearman es la correlación de Pearson de sus rangos:
\[
\rho_S(u,v)
=
\rho_P(\operatorname{rank}(u),\operatorname{rank}(v)).
\]
En 04A se usa para cuantificar asociación monotónica entre distancia en $X$ y
$|y_i-y_j|$ sin asumir relación lineal.

\section{Moran's \texorpdfstring{$I$}{I}}

Con desviaciones $z_i=y_i-\bar y$ y matriz de pesos $W$:
\[
I
=
\frac{n}{S_0}
\frac{\sum_i\sum_jw_{ij}z_i z_j}
{\sum_i z_i^2},
\qquad
S_0=\sum_i\sum_jw_{ij}.
\]
El $p$ reportado en 04A se obtiene por permutación de valores sobre la topología fija. Moran's
$I$ positivo indica que valores similares tienden a ocupar nodos vecinos
\citep{Moran1950}.

\section{Ridge global y local}

Ridge global:
\[
\widehat\beta
=
\arg\min_\beta
\|y-X\beta\|_2^2+\alpha\|\beta\|_2^2.
\]
Si las columnas están centradas y la matriz es de rango apropiado, la solución cerrada es
\[
\widehat\beta
=
(X^\top X+\alpha I)^{-1}X^\top y.
\]

Local Ridge reemplaza la suma global por vecinos y pesos:
\[
\widehat\beta_q
=
\arg\min_\beta
\sum_{i\in\mathcal N_k(q)}
w_{qi}(y_i-x_i^\top\beta)^2
+\alpha\|\beta\|_2^2.
\]
En el finalista,
\[
k=24,\qquad
\alpha=30,\qquad
w_{qi}=(d^{geo}_{qi}+10^{-6})^{-1}.
\]

\section{PLS}

PLS construye componentes latentes supervisados. A nivel conceptual:
\[
t_h=Xw_h,
\]
con direcciones $w_h$ seleccionadas para capturar covarianza con $y$ sujeto a la deflación de
componentes previos \citep{Wold2001}. Dado que $w_h$ depende de $y$, el fit de PLS debe ocurrir
dentro de train en cada fold.

\section{PCA}

PCA es $X$-only. Para matriz centrada $X$, busca
\[
v_1=\arg\max_{\|v\|=1}v^\top X^\top Xv,
\]
y componentes siguientes ortogonales. En 04A, PCA transductivo puede usar las 197 filas porque
no usa etiquetas y la tarea fija permite conocer $X_{\Target}$.

\section{kNN con pesos}

Para $p\ge0$:
\[
w_{qi}
=
(d_{qi}+\epsilon)^{-p},
\qquad
\hat y_q
=
\frac{\sum_{i\in\mathcal N_k(q)}w_{qi}y_i}
{\sum_{i\in\mathcal N_k(q)}w_{qi}}.
\]
$p=0$ corresponde a vecinos con peso uniforme.

\section{Grafo y Laplaciano}

El grafo kNN usa
\[
A_{ij}
=
\exp\left[-\frac12(d_{ij}/\sigma)^2\right]
\]
para aristas seleccionadas y se simetriza tomando el peso máximo en ambas direcciones.
\[
D=\diag(A\mathbf 1),\qquad
L=D-A.
\]
La energía de suavidad:
\[
f^\top Lf
=
\frac12\sum_{ij}A_{ij}(f_i-f_j)^2.
\]
El estimador de 04B/04D resuelve:
\[
(M+\lambda L+\epsilon I)f
=
M(y-\bar y\mathbf1),
\]
donde $M$ tiene uno en nodos etiquetados visibles y cero en nodos sin etiqueta.

\section{Blend convexo}

Para predictores $a$ y $b$:
\[
\hat y^{(w)}=(1-w)\hat y^{(a)}+w\hat y^{(b)},
\qquad w\in[0,1].
\]
Un mínimo aparente de $w$ sobre la misma validación no se toma automáticamente como regla
final. Checkpoint~04F selecciona $w$ leave-one-split-out antes de decidir si la complejidad
adicional está justificada.

\section{Target-matched pseudo-competition}

El pseudo-target count se deriva de:
\[
n_{\Target^\star}
\approx
|\Train|\frac{|\Target|}{|\Train|+|\Target|}
=
138\frac{59}{197}
\approx41.
\]
La construcción de la máscara sólo usa $X$ y metadata. El score del método se calcula
posteriormente con $y_{\Target^\star}$, ya con la predicción congelada.

\section{Discrepancia de finalistas}

En 04F se conserva
\[
u_i=
|\hat y_i^{Local}-\hat y_i^{Graph}|
\]
como medida de sensibilidad a modelo, no como desviación estándar probabilística. Para
interpretarla como intervalo predictivo sería necesaria calibración específica de cobertura,
que no se realizó.
